\chapter{Stochastic Control}\label{ch:qm:stochastic-control}

After a 20\% fall in equities a pension fund's equity weight has drifted from 60\% to 52\%. Its
rulebook says: rebalance to 60\% at the quarter-end, which means buying equities after a fall. Half
the investment committee wants to wait for the market to settle. The question is one of stochastic
control: choose, at each moment and with what is known then, how much to hold, so as to maximise the
expected utility of what the fund ends with. For an investor with constant \omterm{def:qm:stochastic-control:utility}{relative risk aversion} in a
market of lognormal prices, Merton solved it in 1969: hold a constant fraction, whatever has happened,
and rebalance to it. With this fund's assumptions the fraction is 51\%, holding 60\% instead costs
3.6~basis points of certainty-equivalent return a year, and the committee's debate is about a much
smaller number than the uncertainty in its own expected return. This chapter develops \omterm{def:qm:stochastic-control:dp}{dynamic
programming} in discrete time, the \omterm{def:qm:stochastic-control:hjb}{Hamilton--Jacobi--Bellman equation} in continuous time, the
verification theorem that turns a candidate into a proof, Merton's problem, and a sketch of the
\omterm{def:qm:stochastic-control:viscosity}{viscosity solutions} that handle \omterm{def:qm:stochastic-control:problem}{value functions} with kinks.

\section{Dynamic programming}

\begin{definition}[Stochastic control problem, admissible control, value function]\label{def:qm:stochastic-control:problem}
A \emph{stochastic control problem}\index{stochastic control problem} chooses a process $u_t$ that
influences the dynamics of a state $X$, here $dX_t = \mu(t, X_t, u_t)\,dt + \sigma(t, X_t, u_t)\,dW_t$ or its
discrete-time analogue, to maximise $J(t, x; u) = \E_{t,x}[\int_t^Tf(s, X_s, u_s)\,ds + g(X_T)]$. An
\emph{admissible control}\index{admissible control} is an \omterm{def:qm:probability-at-speed:filtration}{adapted process} with values in a set $\mathcal U$ for
which the state equation has a solution and $J$ is well defined. The \emph{value function}\index{value
function} is $V(t, x) = \sup_uJ(t, x; u)$ over admissible controls.
\end{definition}

\begin{theorem}[Dynamic programming principle]\label{thm:qm:stochastic-control:dpp}
For every $t < s \le T$, $V(t, x) = \sup_u\E_{t,x}\bigl[\int_t^sf(r, X_r, u_r)\,dr + V(s, X_s)\bigr]$.
\end{theorem}

\begin{proof}[Partial proof]
In discrete time with finitely many states and controls: any strategy from $t$ earns its rewards up
to $s$ and then at most $V(s, X_s)$, which gives $\le$; following an optimal (or $\varepsilon$-optimal)
continuation from each state at $s$ after any first part gives $\ge$. In continuous time the argument
needs measurable selection of $\varepsilon$-optimal continuations (Fleming and Soner, 2006).
\end{proof}

\begin{definition}[Dynamic programming, feedback control]\label{def:qm:stochastic-control:dp}
\emph{Dynamic programming}\index{dynamic programming} computes the \omterm{def:qm:stochastic-control:problem}{value function} backward in time
with the one-step principle, $V_t(x) = \max_u\{f(t, x, u) + \E[V_{t+1}(X_{t+1}) \mid X_t = x, u]\}$, from $V_T =
g$. A \emph{feedback control}\index{feedback control} is a control of the form $u_t = u^*(t, X_t)$, a function of
the current state; dynamic programming produces one.
\end{definition}

The principle reduces a choice over whole strategies to a sequence of one-step choices, at the cost of
computing $V$ on the whole state space: the curse of dimensionality. The running project's solver
(the tutorial's first listing) does it on a grid, with the expectation computed by
quadrature over the next state and linear interpolation between grid points.

\section{The Hamilton--Jacobi--Bellman equation}

\begin{definition}[Hamilton--Jacobi--Bellman equation]\label{def:qm:stochastic-control:hjb}
For the controlled diffusion of \cref{def:qm:stochastic-control:problem}, the \emph{Hamilton--Jacobi--Bellman
equation}\index{Hamilton--Jacobi--Bellman equation} is
\[
\partial_tV(t, x) + \sup_{u\in\mathcal U}\Bigl\{\mu(t, x, u)\partial_xV + \tfrac12\sigma^2(t, x, u)\partial_{xx}V + f(t, x,
u)\Bigr\} = 0, \qquad V(T, x) = g(x).
\]
\end{definition}

It is the \omterm{def:qm:stochastic-control:dp}{dynamic programming} principle over an interval $[t, t + h]$ with $h \to 0$: expand $V(t + h,
X_{t+h})$ by Itô's formula, take expectations, divide by $h$. For each fixed control it is the backward
equation of chapter~4 with the generator $\mathcal L^u$ of the controlled state; the supremum makes it
nonlinear. The heuristic assumes $V$ is smooth, which is what the verification theorem replaces by an
assumption about a candidate.

\section{Verification}

\begin{theorem}[Verification]\label{thm:qm:stochastic-control:verification}
Let $w \in C^{1,2}$ with polynomial growth solve the HJB equation, and suppose the supremum is attained at
$u^*(t, x)$ such that the state equation with the feedback $u^*$ has a solution. Then $w = V$ and $u^*$ is
optimal, provided the stochastic integrals below are true \omterm{def:qm:probability-at-speed:martingale}{martingales}.
\end{theorem}

\begin{proof}
For any admissible $u$, Itô's formula gives $w(T, X_T) = w(t, x) + \int_t^T(\partial_s + \mathcal L^{u_s})w\,ds + \int_t^T\sigma
\partial_xw\,dW_s$. The HJB equation says $(\partial_s + \mathcal L^u)w + f \le 0$, so, taking expectations, $\E[g(X_T) + \int_t^Tf\,ds] \le
w(t, x)$: $J \le w$. With $u = u^*$ the inequality is an equality, so $J(u^*) = w \ge V \ge J(u^*)$.
\end{proof}

\section{Merton's problem}

\begin{definition}[Utility function, relative risk aversion, CRRA utility]\label{def:qm:stochastic-control:utility}
A \emph{utility function}\index{utility function} $U$ is increasing and concave; the investor maximises
$\E[U(W_T)]$. Its \emph{relative risk aversion}\index{relative risk aversion} is $-wU^{\prime\prime}(w)/U'(w)$.
\emph{CRRA utility}\index{CRRA utility} has constant relative risk aversion $\gamma > 0$: $U(w) =
w^{1-\gamma}/(1 - \gamma)$, and $\ln w$ for $\gamma = 1$.
\end{definition}

\begin{definition}[Certainty equivalent]\label{def:qm:stochastic-control:ce}
The \emph{certainty equivalent}\index{certainty equivalent} of a random wealth $W$ is the sure amount $c$
with $U(c) = \E[U(W)]$; over a horizon $T$, the certainty-equivalent return is $\ln(c/W_0)/T$.
\end{definition}

A fund holds a fraction $\pi_t$ of its wealth in a stock with $dS/S = \mu\,dt + \sigma\,dW$ and the rest in the
\omterm{def:qm:girsanov-and-changes-of-numeraire:numeraire}{money-market account} at rate $r$, so $dW_t = W_t(r + \pi_t(\mu - r))\,dt + W_t\pi_t\sigma\,dW_t$.

\begin{theorem}[Merton]\label{thm:qm:stochastic-control:merton}
For \omterm{def:qm:stochastic-control:utility}{CRRA utility} of terminal wealth, the optimal fraction is constant,
\[
\pi^* = \frac{\mu - r}{\gamma\sigma^2},
\]
and the \omterm{def:qm:stochastic-control:problem}{value function} is $V(t, w) = \frac{w^{1-\gamma}}{1-\gamma}e^{(1-\gamma)\rho(T-t)}$, where $\rho = r + \frac{(\mu -
r)^2}{2\gamma\sigma^2}$ is the certainty-equivalent return of the optimal strategy.
\end{theorem}

\begin{proof}
Try $V(t, w) = \frac{w^{1-\gamma}}{1-\gamma}h(t)$ in the HJB equation
\[
\partial_tV + \sup_\pi\bigl\{w(r + \pi(\mu - r))\partial_wV + \tfrac12w^2\pi^2\sigma^2\partial_{ww}V\bigr\} = 0.
\] Since $w\partial_wV = (1 - \gamma)V$ and $w^2\partial_{ww}V = -\gamma(1 - \gamma)V$, the bracket is
$(1 - \gamma)V(r + \pi(\mu - r) - \tfrac12\gamma\pi^2\sigma^2)$, maximised at $\pi^*$ with value $(1 - \gamma)V\rho$. Hence $h' +
(1 - \gamma)\rho h = 0$, $h(T) = 1$. The candidate is smooth and the wealth equation with constant $\pi^*$ is a
\omterm{def:qm:stochastic-differential-equations:gbm}{geometric Brownian motion}, so \cref{thm:qm:stochastic-control:verification} applies.
\end{proof}

\begin{definition}[Merton fraction]\label{def:qm:stochastic-control:fraction}
The \emph{Merton fraction}\index{Merton fraction} is $\pi^* = (\mu - r)/(\gamma\sigma^2)$.
\end{definition}

The optimal fraction does not depend on wealth or on time: after a fall the fund should buy back to it.
Holding a constant $\pi$ earns the certainty-equivalent return $r + \pi(\mu - r) - \tfrac12\gamma\pi^2\sigma^2$
(\cref{fig:qm:stochastic-control:ce}), so a mistaken fraction costs $\tfrac12\gamma\sigma^2(\pi - \pi^*)^2$ a year. With $\mu =
7\%$, $r = 2\%$, $\sigma = 18\%$ and $\gamma = 3$: $\pi^* = 51.4\%$, a certainty-equivalent return of 3.29\%, and 60\%
costs 3.6~basis points a year. With $\gamma = 1$, log utility, $\pi^* = (\mu - r)/\sigma^2$ is the Kelly fraction of One
Quant Book~2, chapter~29, here 154\%: the growth-optimal investor borrows.

\begin{omfigure}
\begin{tikzpicture}
\begin{axis}[width=11cm,height=4.8cm,
  xlabel={fraction in equities $\pi$},ylabel={CE return (\% a year)},
  tick label style={font=\small},label style={font=\small},
  xmin=0,xmax=1.6,ymin=1.5,ymax=3.5,grid=major,grid style={gray!25},
  table/col sep=comma]
\addplot[omDef,very thick] table[x=pi,y=ce]{figdata/methods/09-stochastic-control/ce.csv};
\addplot[black,only marks,mark=*,mark size=1.8pt] coordinates {(0.5144,3.2860) (0.60,3.2504)};
\addplot[gray,dashed] coordinates {(1.5432,1.5) (1.5432,3.5)};
\node[font=\footnotesize,anchor=east] at (axis cs:1.52,2.9) {Kelly ($\gamma = 1$)};
\end{axis}
\end{tikzpicture}

\omcaption{Certainty-equivalent return of a constant fraction in equities for a CRRA investor with $\gamma = 3$
($\mu = 7\%$, $r = 2\%$, $\sigma = 18\%$): a parabola peaking at the \omterm{def:qm:stochastic-control:fraction}{Merton fraction} (left dot, 51.4\%), flat enough that 60\% (right dot) costs only
3.6~bp a year. The Kelly fraction, optimal for $\gamma = 1$, is 154\%. Data: closed form, the chapter's
tutorial.}\label{fig:qm:stochastic-control:ce}
\end{omfigure}

The flatness has a flip side. The fraction is proportional to the expected excess return, which is known
far less well: estimated from ten years of returns with 18\% volatility, $\mu$ has a standard error of $18\%/\sqrt{10}
= 5.7\%$, and every point of $\mu$ moves $\pi^*$ by ten points (\cref{fig:qm:stochastic-control:sensitivity}).
Merton's rule is exact and its input is uncertain; chapter~14 shrinks the input.

\begin{omfigure}
\begin{tikzpicture}
\begin{groupplot}[group style={group size=2 by 1,horizontal sep=1.6cm},
  width=7.6cm,height=4.8cm,tick label style={font=\small},label style={font=\small},
  grid=major,grid style={gray!25},table/col sep=comma]
\nextgroupplot[xlabel={expected excess return $\mu - r$ (\%)},ylabel={Merton fraction},xmin=0,xmax=10,ymin=0,ymax=1.6,
  legend style={font=\footnotesize,at={(0.03,0.97)},anchor=north west,fill=white,draw=none}]
\addplot[omThm,very thick] table[x=excess,y=g2]{figdata/methods/09-stochastic-control/sensitivity.csv};
\addplot[omDef,very thick] table[x=excess,y=g3]{figdata/methods/09-stochastic-control/sensitivity.csv};
\addplot[omProp,very thick] table[x=excess,y=g5]{figdata/methods/09-stochastic-control/sensitivity.csv};
\legend{$\gamma = 2$,$\gamma = 3$,$\gamma = 5$}
\nextgroupplot[xlabel={years},ylabel={equity weight},xmin=0,xmax=10,ymin=0.4,ymax=0.9]
\addplot[name path=hi,gray!40,forget plot] table[x=year,y=p90]{figdata/methods/09-stochastic-control/drift.csv};
\addplot[name path=lo,gray!40,forget plot] table[x=year,y=p10]{figdata/methods/09-stochastic-control/drift.csv};
\addplot[gray!20,forget plot] fill between[of=hi and lo];
\addplot[omDef,very thick] table[x=year,y=p50]{figdata/methods/09-stochastic-control/drift.csv};
\addplot[black,dashed,domain=0:10] {0.5144};
\end{groupplot}
\end{tikzpicture}

\omcaption{Left: the \omterm{def:qm:stochastic-control:fraction}{Merton fraction} against the expected excess return, for three risk aversions ($\sigma = 18\%$):
ten points of fraction per point of return at $\gamma = 3$. Right: the equity weight of a fund that starts at 60\% and
never rebalances, median and 10th--90th percentile band over 20\,000 simulated paths, against the \omterm{def:qm:stochastic-control:fraction}{Merton fraction}
(dashed). Data: the chapter's tutorial, seeded.}\label{fig:qm:stochastic-control:sensitivity}
\end{omfigure}

\section{A sketch of viscosity solutions}

\omterm{def:qm:stochastic-control:problem}{Value functions} are often not smooth: a control that switches from one extreme to another at a boundary puts a kink
in $V$, and the HJB equation cannot hold there in the classical sense. The right notion replaces derivatives by
smooth test functions touching $V$.

\begin{definition}[Viscosity solution]\label{def:qm:stochastic-control:viscosity}
A continuous $V$ is a \emph{viscosity solution}\index{viscosity solution} of $F(x, V, DV, D^2V) = 0$ ($F$ nonincreasing
in its last argument) if at every point where a smooth $\phi$ touches $V$ from above ($V - \phi$ has a local maximum),
$F(x, V, D\phi, D^2\phi) \le 0$, and at every point where it touches from below, $F \ge 0$.
\end{definition}

For the problem of leaving $(-1, 1)$ as fast as possible at unit speed, $V(x) = 1 - |x|$ and $|V'| = 1$ everywhere but at
zero, where no derivative exists; it is the \omterm{def:qm:stochastic-control:viscosity}{viscosity solution}, and it is the limit as $\varepsilon \to 0$ of the
smooth solutions of $-\varepsilon u^{\prime\prime} + |u'| = 1$ (\cref{fig:qm:stochastic-control:viscosity}), which is where the
name comes from. Crandall and Lions (1983) proved existence and uniqueness for first-order equations; the theory
guarantees that \omterm{def:qm:stochastic-control:dp}{dynamic programming} on a grid converges to the right \omterm{def:qm:stochastic-control:problem}{value function} even when it has kinks, as in
the stopping problems of chapter~10.

\begin{omfigure}
\begin{tikzpicture}
\begin{axis}[width=10cm,height=4.6cm,
  xlabel={$x$},ylabel={$u(x)$},
  tick label style={font=\small},label style={font=\small},
  xmin=-1,xmax=1,ymin=0,ymax=1.1,grid=major,grid style={gray!25},
  legend columns=3,legend style={font=\footnotesize,at={(0.5,-0.3)},anchor=north,draw=none,fill=none,/tikz/every even column/.append style={column sep=8pt}},
  table/col sep=comma]
\addplot[omProp,very thick] table[x=x,y=e02]{figdata/methods/09-stochastic-control/viscosity.csv};
\addplot[omThm,very thick] table[x=x,y=e005]{figdata/methods/09-stochastic-control/viscosity.csv};
\addplot[omDef,very thick,dashed] table[x=x,y=e0]{figdata/methods/09-stochastic-control/viscosity.csv};
\legend{$\varepsilon = 0.2$,$\varepsilon = 0.05$,$1 - |x|$}
\end{axis}
\end{tikzpicture}

\omcaption{Solutions of $-\varepsilon u^{\prime\prime} + |u'| = 1$ on $(-1, 1)$ with $u(\pm1) = 0$, namely $1 - |x| + \varepsilon(e^{-1/\varepsilon} -
e^{-|x|/\varepsilon})$: smooth for every $\varepsilon > 0$, converging to the kinked \omterm{def:qm:stochastic-control:problem}{value function} $1 - |x|$, the \omterm{def:qm:stochastic-control:viscosity}{viscosity
solution} of $|u'| = 1$.}\label{fig:qm:stochastic-control:viscosity}
\end{omfigure}

\section{Tutorial: Merton by backward induction}\label{tut:qm:stochastic-control:merton}

\begin{tutorial}
\textbf{Goal.} Solve a discrete-time version of the fund's problem by \omterm{def:qm:stochastic-control:dp}{dynamic programming}, check it against
Merton's fraction, and measure what rebalancing buys. \textbf{End state:}
\cref{fig:qm:stochastic-control:ce,fig:qm:stochastic-control:sensitivity} and the numbers of the weekend problem.
\begin{tutsteps}
\item \textbf{The solver}: backward induction on a grid, with an optional stopping reward for chapter~10.
\omcode{code/firm/dpsolve/firm_dpsolve.py}{30}{49}{Backward induction with quadrature and interpolation.}\label{lst:qm:stochastic-control:dp}
\item \textbf{The problem}: ten annual rebalancing dates, log wealth on a grid of 241 points, fractions from 0 to 1.5
by steps of 0.005, lognormal annual returns by 20-point Gauss--Hermite quadrature.
\omcode{code/methods/09-stochastic-control/python/qm_merton.py}{66}{84}{Merton's problem as a grid dynamic programme.}\label{lst:qm:stochastic-control:merton}
\item \textbf{Run} \texttt{problem()}: the optimal fraction is 0.52 at every wealth level and date, the grid point
nearest the one-period optimum 0.517 (continuous rebalancing gives 0.514); then \texttt{fig\_merton.py}.
\end{tutsteps}
\textbf{What to change next.} Add consumption at a constant fraction of wealth and check Merton's consumption rule;
replace CRRA by a utility with a floor (a liability the fund must cover) and watch the fraction start to depend on
wealth.
\end{tutorial}

\section{Build: the dynamic-programming solver}\label{bld:qm:stochastic-control:dpsolve}

\begin{build}
\bfield{Purpose} Solve the firm's small control problems (allocation, inventory, execution schedules) and, with a
stopping reward, its stopping problems (chapter~10), by backward induction on a grid.
\bfield{Interface} \texttt{backward\_induction(grid, controls, n\_steps, reward, transition, terminal, stop=None)}
returning values, policy and the stopping region; \texttt{gauss\_hermite\_normal(n)}.
\bfield{Rules} One-dimensional state grid, finite control set, expectations by user-supplied nodes and weights, linear
interpolation, flat extrapolation; vectorised over states and controls.
\bfield{Acceptance tests} \texttt{code/firm/dpsolve/tests/}: Gauss--Hermite integrates polynomials exactly; the Merton
problem returns a wealth-independent fraction equal to the one-period optimum; a stopping problem with a known
solution (the American-style exit of chapter~10) is reproduced.
\bfield{Stretch} Two-dimensional states with bilinear interpolation; policy iteration for infinite horizons.
\end{build}

\begin{omsources}
\item R.~C.~Merton, ``Lifetime portfolio selection under uncertainty: the continuous-time case'', \emph{Review of
Economics and Statistics} 51, 1969.
\item P.~A.~Samuelson, ``Lifetime portfolio selection by dynamic stochastic programming'', \emph{Review of Economics and
Statistics} 51, 1969.
\item R.~Bellman, \emph{Dynamic Programming}, Princeton University Press, 1957.
\item M.~G.~Crandall and P.-L.~Lions, ``Viscosity solutions of Hamilton--Jacobi equations'', \emph{Transactions of the
AMS} 277, 1983.
\item W.~H.~Fleming and H.~M.~Soner, \emph{Controlled Markov Processes and Viscosity Solutions}, Springer, 2nd ed., 2006.
\end{omsources}

\section{Exercises}

\begin{exercise}[$\star$]\label{exo:qm:stochastic-control:1}
An asset has an expected excess return of 4\% and a volatility of 20\%. What fraction does an investor with $\gamma = 4$
hold?
\end{exercise}

\begin{exercise}[$\star$]\label{exo:qm:stochastic-control:2}
With the fund's assumptions, what is the certainty-equivalent return of holding 100\% in equities?
\end{exercise}

\begin{exercise}[$\star$]\label{exo:qm:stochastic-control:3}
Compute the \omterm{def:qm:stochastic-control:utility}{relative risk aversion} of $U(w) = \ln w$ and of $U(w) = w^{1-\gamma}/(1 - \gamma)$.
\end{exercise}

\begin{exercise}[$\star\star$]\label{exo:qm:stochastic-control:4}
Show that holding a constant $\pi$ instead of $\pi^*$ costs $\tfrac12\gamma\sigma^2(\pi - \pi^*)^2$ of certainty-equivalent return a year,
and evaluate it for 60\% against 51.4\%.
\end{exercise}

\begin{exercise}[$\star\star$]\label{exo:qm:stochastic-control:5}
What fraction does a log-utility investor hold with the fund's assumptions? Why would a fund not do so?
\end{exercise}

\begin{exercise}[$\star\star$]\label{exo:qm:stochastic-control:6}
Verify that $V = \frac{w^{1-\gamma}}{1-\gamma}e^{(1-\gamma)\rho(T-t)}$ solves the HJB equation, with $\rho = r + (\mu - r)^2/(2\gamma\sigma^2)$, and
compute $\rho$ for the fund.
\end{exercise}

\begin{exercise}[$\star\star\star$]\label{exo:qm:stochastic-control:7}
\emph{Coding.} With \texttt{backward\_induction}, solve the fund's problem with $\gamma = 5$ and annual rebalancing. Is
the fraction still independent of wealth, and how does it compare with $(\mu - r)/(\gamma\sigma^2)$?
\end{exercise}

\begin{exercise}[$\star\star\star$]\label{exo:qm:stochastic-control:8}
\emph{Find the flaw.} ``We estimated the equity premium at 5\% from ten years of data, so our optimal equity weight is
51\%, and we will hold exactly that.''
\end{exercise}

\section{Problem: The Rebalancing Committee}

\begin{problem}[{Weekend problem --- how much the committee's decision is worth}]\label{pb:qm:stochastic-control:1}
A fund with \omterm{def:qm:stochastic-control:utility}{CRRA utility}, $\gamma = 3$, expects equities to return $\mu = 7\%$ with volatility $\sigma = 18\%$; cash earns $r =
2\%$. Its rulebook holds 60\% in equities, rebalanced. After a fall its weight is 52\%.

\medskip\noindent\textbf{Part I --- Merton.}
\begin{enumerate}
\item What is the \omterm{def:qm:stochastic-control:fraction}{Merton fraction}?
\item What is the certainty-equivalent return of the optimal strategy?
\item What does holding 60\% cost a year?
\item And holding 52\%?
\item What trade does Merton recommend after the fall, and what does the rulebook recommend?
\end{enumerate}

\medskip\noindent\textbf{Part II --- Sensitivity.}
\begin{enumerate}[resume]
\item What are the fractions for $\gamma = 2$ and $\gamma = 5$?
\item And for $\mu = 6\%$ and $8\%$?
\item What does $\gamma = 1$ give, and what does it mean?
\item What does the grid dynamic programme with annual rebalancing give, and the one-period optimum?
\item Why does the fraction not depend on wealth?
\end{enumerate}

\medskip\noindent\textbf{Part III --- Rebalancing.}
\begin{enumerate}[resume]
\item Over ten simulated years, what certainty-equivalent returns do constant 60\%, Merton, and buy-and-hold from 60\% achieve?
\item Where does the buy-and-hold weight end after ten years (10th percentile, median, 90th)?
\item What does never rebalancing cost, against rebalancing to 60\%?
\item With what standard error is $\mu$ known from ten years of returns, and what band does that put on $\pi^*$?
\item Which matters more for the fund: the committee's choice between 52\% and 60\%, or the uncertainty in $\mu$?
\end{enumerate}

\medskip\noindent\textbf{Part IV --- Judgement.}
\begin{enumerate}[resume]
\item What would transaction costs change?
\item What would liabilities (a pension promise) change?
\item What should the chief investment officer tell the committee?
\item State the \emph{named result}: the \omterm{def:qm:stochastic-control:fraction}{Merton fraction} and the cost of the rulebook's 60\%.
\item In one sentence: what does Merton's solution say about rebalancing after a fall?
\end{enumerate}
\end{problem}

\section{Interview questions}

\begin{interviewq}[$\star$ \iqroles{researcher, trader}]\label{iq:qm:stochastic-control:1}
State Merton's optimal fraction and explain each of its three inputs.
\end{interviewq}

\begin{interviewq}[$\star\star$ \iqroles{researcher}]\label{iq:qm:stochastic-control:2}
Derive the HJB equation from the \omterm{def:qm:stochastic-control:dp}{dynamic programming} principle.
\end{interviewq}

\begin{interviewq}[$\star\star$ \iqroles{researcher}]\label{iq:qm:stochastic-control:3}
Why does a CRRA investor in Merton's model buy equities after they fall?
\end{interviewq}

\begin{interviewq}[$\star\star$ \iqroles{trader, researcher}]\label{iq:qm:stochastic-control:4}
How does the Kelly criterion relate to Merton's problem?
\end{interviewq}

\begin{interviewq}[$\star\star$ \iqroles{developer}]\label{iq:qm:stochastic-control:5}
Your dynamic programme on a grid returns a policy that jumps between neighbouring grid points. What do you check?
\end{interviewq}

\begin{interviewq}[$\star\star\star$ \iqroles{researcher}]\label{iq:qm:stochastic-control:6}
When is the \omterm{def:qm:stochastic-control:problem}{value function} of a control problem not smooth, and how does one make sense of the HJB equation then?
\end{interviewq}
